\begin{helper}[Transport through the two switches]
Let $X=\noderef{ScheideEnlargedVertexSet}(R)$ and suppose that $X$ is not
strongly closed.  Then there are a color $\alpha$, a defective index
$d=(\delta,h_\delta)$, an edge $f$, and vertices $w,u$ with the following
properties.  Put
\[
 \varphi_d=\noderef{TwoColorComponentSwap}
 (G,\varphi,\gamma(d),\delta,\operatorname{inl}v^0_d),
\]
\[
 Q=\noderef{TwoColorComponent}
 (G,\varphi_d,\alpha_0,\alpha,\operatorname{inl}u),\qquad
 \psi=\noderef{TwoColorComponentSwap}
 (G,\varphi_d,\alpha_0,\alpha,\operatorname{inl}u).
\]
Then $\alpha\ne\alpha_0$, $G.\mathop{\rm ends}(f)=s(w,u)$,
$\varphi_d(f)=\alpha$, $\operatorname{inr}f\in Q$, and $Q$ contains no
vertex-node from $V(T)$.  Furthermore
\[
 (T_d,e_0,\psi),\qquad T_d=(R d).\mathit{replacementTree},
\]
is a globally maximum triple; $\alpha_0$ is missing at $y_0=C.y(0)$
under $\varphi_d$; $y_0,w\in V(T_d)$; $u\notin V(T_d)$; and $f$ is not
an edge identity of $T_d$.
\end{helper}
\begin{proof}
Apply \noderef{ScheideEnlargedBoundarySelection}.  It supplies
$\alpha,d,f,w,u$ with $\alpha\ne\alpha_0$, the stated endpoint equation,
$\varphi(f)=\alpha$, membership of $f$ in the original component through
$u$, avoidance of $V(T)$ by that component, and
$w=(R d).u_1$ or $w=(R d).u_2$.  Write $\delta=d.1$,
$r=D.\gammaIndex(d)$, $\gamma=C.\gamma(r)$, and
$T_d=(R d).\mathit{replacementTree}$.

We first separate the four colors relevant to the first switch.  According
to which of $(R d).u_1,(R d).u_2$ equals $w$, the corresponding one of
$(R d).f_1,(R d).f_2$ is a $\delta$-edge incident with $w$ whose other
end is in the retained part of $V(T)$.  The second clause of
\noderef{ScheideSymmetricDefectiveInternal} supplies a $\gamma$-edge
internal to the pair $\{(R d).u_1,(R d).u_2\}$ and hence incident with
$w$.  These identities differ from $f$, because both of their ends lie
in $X$, whereas $f$ leaves $X$.  Properness from the
\noderef{GlobalMaximumTriple} inside
\noderef{ScheideStructuralConfiguration}, together with
\noderef{AdjacentEdges}, gives
$\alpha\ne\delta$ and $\alpha\ne\gamma$.
The setup gives $\alpha_0\ne\gamma$.  Since $\delta$ is defective on
$V(T)$, a $\delta$-edge leaves $V(T)$; closedness and the missingness of
$\alpha_0$ at $y_0$ exclude equality $\delta=\alpha_0$.  Thus neither
$\alpha$ nor $\alpha_0$ is changed into the other by the
$(\gamma,\delta)$ switch.

Unfold \noderef{TwoColorComponentSwap}.  For every edge identity $q$,
the preceding inequalities give both equivalences
\[
 \varphi_d(q)=\alpha_0\iff\varphi(q)=\alpha_0,
 \qquad
 \varphi_d(q)=\alpha\iff\varphi(q)=\alpha.                 \tag{1}
\]
Indeed the reverse implications are fixed branches of the switch, and a
new value in either forward implication could only have arisen from
$\gamma$ or $\delta$, which the inequalities exclude.  Hence the labeled
$(\alpha_0,\alpha)$ graphs in \noderef{TwoColorIncidenceGraph} are equal.
The component $Q$ under $\varphi_d$ is therefore the component furnished
under $\varphi$ by \noderef{ScheideEnlargedBoundarySelection}.  It
contains $\operatorname{inr}f$, avoids all vertex-nodes of $V(T)$, and
(1) gives $\varphi_d(f)=\alpha$.

No edge identity of $T_d$ belongs to $Q$.  By the exact edge equation in
\noderef{ScheideReplacementData}, such an identity is either a retained
edge of $T$ or one of $(R d).f_1,(R d).f_2$.  Every old tree edge is
incident with a vertex of $V(T)$ by \noderef{TreeSequence}, and an edge
of $Q$ incident with such a vertex would put that vertex-node in $Q$.
The two new identities have old color $\delta$; equations (1) and
$\delta\notin\{\alpha_0,\alpha\}$ exclude them from the labeled
component.  In particular $f$ is not an edge identity of $T_d$.

The theorem \noderef{ScheideReplacementMaximum} says that
$(T_d,e_0,\varphi_d)$ is globally maximum.  We now verify that its second
switch still has all Tashkinov witnesses.  For a retained noninitial base
edge, the original \noderef{TashkinovTree} condition supplies an earlier
missing witness in the base.  Its color differs from every block color by
the base-edge avoidance clause of \noderef{ScheideStructuralConfiguration}.
For a retained block edge, its color is the corresponding block color,
and the structural configuration supplies a base vertex missing that
color.  Injectivity of the block-color list shows that, outside the
deleted block, this color differs from $\gamma$.  In either case the
position equations of \noderef{ScheideReplacementData} retain the
witness before the retained or shifted edge.

The exact effect in \noderef{ScheideDefectiveSwitchEffect} fixes every
retained tree edge (it is incident with $V(T)$ and is not $D.f_0(d)$).
It fixes the missing set at every retained witness except possibly
$v^0_d$; at $v^0_d$ it erases only $\gamma$ and inserts $\delta$, so it
still retains every witness color just considered.  Every base vertex is
in the retained set because the base and block positions are distinct by
vertex injectivity of \noderef{TreeSequence}.

For the final two edges of $T_d$, the replacement data give color
$\delta$, and \noderef{ScheideDefectiveSwitchEffect} inserts $\delta$
into the missing set at $v^0_d$.  To locate that witness in the retained
base, unpack \noderef{ScheideDefectiveSetup}: its recorded full simple
$(\alpha_0,\delta)$ path starts at $y_0$, ends outside $V(T)$, contains
$v^0_d$, and has no later vertex-node in $V(T)$.  Reversing the displayed
order of the two colors does not change \noderef{TwoColorIncidenceGraph},
so these data give the predicate \noderef{ScheideExitVertex} for
$v^0_d$.  Thus $v^0_d$ belongs to \noderef{ScheideExitVertices}, and
\noderef{ScheideStructuralExitsInBase} writes $v^0_d=C.y(a)$ for some
$a<C.p$.  The deleted vertices occur at positions $C.p+2r$ and
$C.p+2r+1$, so vertex injectivity retains $v^0_d$.  The replacement
position equations retain it at position $a$, and
$a\le C.p-1<T.length-1$, before both final edges.  It is therefore the
required missing-$\delta$ witness for each.

All missing witnesses just constructed lie in $V(T)$.  The component
$Q$ avoids every vertex-node there.  By the incidence definition, no
edge of $Q$ is incident with a witness, so the second switch preserves
each witness's missing set.  It also fixes every tree-edge identity and
color because no edge identity of $T_d$ belongs to $Q$.  Consequently
$T_d$ remains a Tashkinov tree under $\psi$.

Since $\alpha\ne\alpha_0$, properness is preserved by
\noderef{TwoColorComponentSwapProper}.  The switch fixes the value
``uncolored'' and permutes only two colored values, so the exactly-one-
uncolored clause in \noderef{ScheideReplacementMaximum} is preserved.
The criticality, chromatic-index, and degree clauses of
\noderef{KTriple} depend only on $G$ and $k$.  Hence $(G,e_0,\psi)$ is a
$k$-triple.  The tree and its vertex set did not change, so the maximum
cardinality equation in \noderef{ScheideReplacementMaximum}, together
with the Tashkinov verification, proves that $(T_d,e_0,\psi)$ is globally
maximum.

Finally, $y_0$ is a retained base vertex of $T_d$.  Under the first
switch, \noderef{ScheideDefectiveSwitchEffect} either fixes its missing
set or, if $y_0=v^0_d$, erases $\gamma\ne\alpha_0$ and inserts $\delta$.
Thus $\alpha_0$ is missing at $y_0$ under $\varphi_d$.  The replacement
vertex equation places the selected $w=(R d).u_1$ or $(R d).u_2$ in
$V(T_d)$ and gives
\[
 V(T_d)\subseteq V(T)\cup\{(R d).u_1,(R d).u_2\}\subseteq X.
\]
Since $u\notin X$, it is not in $V(T_d)$.  These are all the asserted
transport properties.
\end{proof}
